\section{Meta-Reinforcement Learning con MAML y Reptile}

\subsection*{a) Predicciones, pérdida y gradiente inicial}

Con $(w,b)=(1,0)$ se obtiene $f_\theta(1)=1$ y $f_\theta(2)=2$.
Los residuos $e_j=wx_j+b-y_j$ son $e_1=-2$ y $e_2=-3$, por lo que
\[
    L_{\mathcal{T}_1}(\theta)
    =\frac{(-2)^2+(-3)^2}{2}=6.5.
\]
Al derivar el error cuadrático medio sobre las dos muestras,
\begin{align*}
    \frac{\partial L_{\mathcal{T}_1}}{\partial w}
    &=\frac{2}{K}\sum_{j=1}^K e_j x_j
      =(-2)(1)+(-3)(2)=-8,\\
    \frac{\partial L_{\mathcal{T}_1}}{\partial b}
    &=\frac{2}{K}\sum_{j=1}^K e_j
      =-2-3=-5.
\end{align*}

\subsection*{b) Un paso de adaptación interna}

Con $\alpha=0.1$, se actualiza la inicialización de MAML para la
tarea $\mathcal{T}_1$~\cite{finn2017},
\[
    \theta'_1=(1,0)^\top-0.1(-8,-5)^\top
    =\boxed{(1.8,\,0.5)^\top}.
\]
El modelo adaptado predice $f_{\theta'_1}(1)=2.3$ y
$f_{\theta'_1}(2)=4.1$. En el mismo conjunto de soporte,
\[
    L_{\mathcal{T}_1}(\theta'_1)
    =\frac{(2.3-3)^2+(4.1-5)^2}{2}
    =\frac{0.49+0.81}{2}=\boxed{0.65}.
\]
El paso redujo en un 90 por ciento la pérdida sobre las muestras usadas
para adaptar el modelo.

\subsection*{c) Metagradiente, FOMAML y Reptile}

Sea $L_i^{\mathrm{soporte}}$ la pérdida usada para la adaptación en
la tarea $i$. Sea $L_i^{\mathrm{consulta}}$ la pérdida externa,
evaluada tras la adaptación y generalmente sobre muestras
distintas de la misma tarea. Para un paso interno,
\[
    \theta_i'=\theta-\alpha\nabla_\theta L_i^{\mathrm{soporte}}(\theta),
    \qquad
    \frac{\partial\theta_i'}{\partial\theta}
    =I-\alpha\nabla_\theta^2 L_i^{\mathrm{soporte}}(\theta).
\]
Con gradientes como vectores columna, la regla de la cadena da
\[
    \nabla_\theta L_i^{\mathrm{consulta}}(\theta_i')
    =\Bigl(I-\alpha\nabla_\theta^2 L_i^{\mathrm{soporte}}(\theta)\Bigr)^\top
      \nabla_{\theta_i'}L_i^{\mathrm{consulta}}(\theta_i').
\]
La Hessiana aparece porque el gradiente del paso interno depende
de $\theta$. En el regresor del enunciado, la Hessiana del
soporte es $\begin{pmatrix}5&3\\3&2\end{pmatrix}$, obtenida al
derivar de nuevo $\frac{1}{2}[(w+b-3)^2+(2w+b-5)^2]$. MAML
puede evaluar la contribución de la Hessiana mediante un producto
Hessiana-vector, sin formar la matriz completa~\cite{finn2017}.

\textbf{FOMAML.}
Aproxima el jacobiano de la adaptación por $I$. Usa
$\nabla_{\theta_i'}L_i^{\mathrm{consulta}}(\theta_i')$
para actualizar la meta-inicialización. Conserva el cálculo
del gradiente sobre las muestras de consulta en los parámetros
adaptados y omite la derivación a través de la
adaptación~\cite{finn2017,nichol2018}.

\textbf{Reptile.}
Aplica uno o más pasos de gradiente de la tarea y desplaza la
meta-inicialización hacia los parámetros resultantes
mediante $\theta\leftarrow\theta+\epsilon(\theta_i'-\theta)$.
Tras un único paso interno,
\[
    \theta+\epsilon(\theta_i'-\theta)
    =\theta-\epsilon\alpha\nabla_\theta L_i^{\mathrm{soporte}}(\theta).
\]
En ese caso, Reptile equivale al descenso de gradiente
sobre la pérdida de soporte de la tarea muestreada. Con
varios pasos internos sobre distintos lotes de una tarea,
el desplazamiento incorpora el efecto de la adaptación
entre lotes~\cite{nichol2018}.

Para $d$ parámetros, la combinación de vectores de la
actualización externa cuesta $O(d)$ en ambos métodos
una vez obtenidos el gradiente de consulta en FOMAML o
los parámetros adaptados en Reptile. El costo total incluye
evaluar los gradientes de primer orden y depende también del
modelo, los datos y el número de pasos internos.
